\documentclass[12pt,a4paper,oneside]{amsart}
\usepackage{fullpage,url,amssymb,enumerate,colonequals}
% add a bit of space between text and page numbers at the bottom of the page:
\addtolength{\footskip}{10pt}

\usepackage[all,pdf]{xy} % for complicated commutative diagrams, use pdf driver

\usepackage{comment}
\usepackage{color}
\usepackage{charter,euler} % for better looks
% Make sure that math operators are set in Charter, too.
\DeclareSymbolFont{bchoperators}{T1}{bch}{m}{n}
\SetSymbolFont{bchoperators}{bold}{T1}{bch}{b}{n}
\makeatletter
\renewcommand{\operator@font}{\mathgroup\symbchoperators}
\makeatother

\DeclareFontEncoding{OT2}{}{} % to enable usage of cyrillic fonts
\newcommand{\textcyr}[1]{%
 {\fontencoding{OT2}\fontfamily{wncyr}\fontseries{m}\fontshape{n}\selectfont #1}}
\newcommand{\Sha}{{\mbox{\textcyr{Sh}}}}

\usepackage{titlesec} % change small caps in section titles to boldface
\titleformat{\section}{\normalfont\bfseries\filcenter}{\thesection}{1em}{}
\titleformat{\subsection}{\normalfont\bfseries}{\thesubsection}{1em}{}
\titleformat{\subsubsection}{\normalfont\itshape}{\thesubsubsection}{1em}{}

\definecolor{darkgreen}{rgb}{0,0.5,0}

\usepackage[
        colorlinks, citecolor=darkgreen,
        pdfauthor={Michael Stoll},
        pdftitle={Documentation for the ratpoints program},
]{hyperref}

\newcommand{\Q}{{\mathbb Q}}
\newcommand{\R}{{\mathbb R}}
\newcommand{\BP}{{\mathbb P}}

\newcommand{\rpversion}{2.2.3} % don't forget the verbatim stuff!

\addtolength{\hoffset}{-1cm}
\addtolength{\textwidth}{2cm}

\setlength{\parskip}{0.8ex plus 0.1ex minus 0.1ex}
\setlength{\parindent}{0mm}

%%%%%%%%%%%%%%%%%%%%%%%%%%%%%%%%%%%%%%%%%%%%%%%%%%%%%%%%%%%%%%%%%%%%%%%%

\begin{document}

\title{Documentation for the ratpoints program}

\author{Michael Stoll}
\address{Mathematisches Institut,
         Universit\"at Bayreuth,
         95440 Bayreuth, Germany.}
\email{Michael.Stoll@uni-bayreuth.de}
\date{May 15, 2023}

\maketitle

%=========================================================================

\section{Introduction}

This paper describes the \texttt{ratpoints} program. This program tries to
find all rational points within a given height bound on a hyperelliptic
curve in the most efficient way possible.

%=========================================================================

\section{History and Acknowledgments}

This program goes back to an implementation of the `quadratic sieving'
idea by \textbf{Noam Elkies} that was around in the early 1990s. My own
first contribution was to replace the \texttt{char} arrays that were used
to store the sieving information by bit arrays, in~1995. \textbf{Colin Stahlke}
then made use of the \texttt{gmp} library, so that points could be checked
exactly, and implemented the selection of sieving primes according to
their likely success rate, in~1998. After that, I successively put in
numerous improvements (and some bug fixes). For details of how the
program works, see Section~\ref{impl} below.

Along with \textbf{Noam Elkies} and \textbf{Colin Stahlke}, I would like to
thank \textbf{John Cremona} and \textbf{Sophie Labour} for bug reports and
suggestions for improvements. Further thanks are due to \textbf{Bill Allombert}
for a first implementation of the use of 256-bit (and possibly 512-bit) registers.

%=========================================================================

\section{Availability}

The \texttt{ratpoints-\rpversion} package can be downloaded from my homepage,
see~\cite{ratpoints}, or from \texttt{github} at
\url{https://github.com/MichaelStollBayreuth/ratpoints}.

This program is free software: you can redistribute it and/or modify
it under the terms of the GNU General Public License as published
by the Free Software Foundation, either version~2 of the License, or
(at your option) any later version.

This program is distributed in the hope that it will be useful,
but \textbf{without any warranty}; without even the implied warranty of
\textbf{merchantability} or \textbf{fitness for a particular purpose}. See the
GNU General Public License for more details.

You should receive a copy of the GNU General Public License
% and the GNU Lesser General Public License
along with this program.

If not, see \url{http://www.gnu.org/licenses/}.

%=========================================================================

\section{Installation}

This section describes the installation procedure under Linux.

\subsection{Extract the archive}

\begin{verbatim}
> tar xzf ratpoints-2.2.3.tar.gz
\end{verbatim}

This sets up a directory \texttt{ratpoints-\rpversion} containing the various
files that belong to the installation.

\subsection{Build the program and library}

Do

\begin{verbatim}
> cd ratpoints-2.2.3
> make all
\end{verbatim}

This will build the library \texttt{libratpoints.a}, the executable
\texttt{ratpoints}, and also run pdf\/latex to generate this documentation.

Two of the constants that govern the sieve depend on the machine rather than
on the curve. Reasonable values are compiled in, but \texttt{make tune} will
measure better ones for your machine if you want it to; see
Section~\ref{sec:tuning}. It is a separate step because it takes several
minutes.

By default, the program has primes up to 251 available, and considers the
smallest $30$ of them; for a curve that needs more, it will look further of
its own accord, as described in Section~\ref{sec:spchoice}. If you intend to
do computations with large height bounds, it may make sense to increase the
range of primes further. This can be achieved via

\begin{verbatim}
> make all PRIME_SIZE=8
\end{verbatim}

perhaps after a

\begin{verbatim}
> make distclean
\end{verbatim}

to remove the files that were generated previously. The \texttt{PRIME\_SIZE}
argument can be given any value from 5 to~10; otherwise it is taken to be~7.
The precise meaning is that the program will work with primes $< 2^s$,
when $\text{\texttt{PRIME\_SIZE}} = s$.

The sieve works on several numerators at a time, in a register whose width is
selected by the line
\begin{verbatim}
CCFLAGS1 = ${CCFLAGS256}
\end{verbatim}
in \texttt{Makefile}. It can be set to \verb+${CCFLAGS64}+ (64-bit
registers, i.e., plain \texttt{unsigned long}s, which works on any machine
the library builds on at all), \verb+${CCFLAGS128}+,
\verb+${CCFLAGS256}+ (the default, using 256-bit AVX registers) or
\verb+${CCFLAGS512}+. The \texttt{Makefile} notices that the setting has
changed and rebuilds everything that depends on it, so a \texttt{make
distclean} in between is not necessary. Wider
registers help the first phase of the sieve, where most of the time goes on a
curve with few rational points, but they help the second phase much less and
the sieve tables not at all, so the gain is nowhere near proportional to the
width: on the author's machine, 64, 128 and 256~bits take 21.3, 15.2
and~12.2 units of time on such a curve.

The 128-bit setting uses the macro \verb+USE_AVX128+. In spite of its name it
needs only SSE2 instructions, so it runs wherever the older \verb+USE_SSE+
variant did, and it is never slower, because its test for an empty register
compares the whole register against zero instead of extracting both halves
into general registers. That is worth about 30\% of the second phase of the
sieve, which comes to 9\% of a run long enough for the sieve to dominate it
and about 3\% of \texttt{make test1}. The older variant is still available as
\verb+${CCFLAGS128s}+.

You may have to change \verb+-mavx2+ in the line
\begin{verbatim}
CCFLAGS256 = -DUSE_AVX -mavx2
\end{verbatim}
into \verb+-mavx+ or leave it out altogether, depending on the capabilities
of your CPU. You will notice that AVX2 (or AVX) is not supported, when the
executable throws an `unknown instruction' error or similar.

There is also a variant using 512-bit registers, activated by
\begin{verbatim}
CCFLAGS1 = ${CCFLAGS512}
\end{verbatim}
which requires a CPU with AVX512F capability. Note that this variant has
so far only been tested indirectly (see the change log entry below), since
I have no such CPU available; you should therefore run \texttt{make test}
before relying on it. It is also worth timing the 256-bit and the 512-bit
version against each other on your machine: the sieving loop is limited by
how fast it loads bit arrays from the first-level cache and by how much of
that cache the sieve tables occupy, rather than by arithmetic, so wider
registers do not automatically pay off, and some Intel server processors reduce their clock
frequency when 512-bit instructions are used. The one data point available
so far is encouraging, though: on an AMD Ryzen~9 9950X3D (Zen~5 architecture),
\textbf{Drew Sutherland} measured a speed-up of 13 to 14\% over the 256-bit
version, for a curve with many rational points.

If you leave out \verb+-mavx512f+ and compile with just
\verb+-DUSE_AVX512+, then \texttt{gcc} will express the 64-byte vectors
through narrower operations, so that the code runs on any machine. This is
of no use in production (it is slower than the 256-bit version), but it
does allow the 512-bit code path to be tested without the corresponding
hardware. \texttt{gcc} will warn that passing a 512-bit vector argument
without AVX512F enabled changes the ABI; this is harmless, as the whole
library is built with the same flags.

\subsection{Run a test}

Run

\begin{verbatim}
> make test1
\end{verbatim}

in the working directory. This will build an executable \texttt{rptest} and
then run (and time) it. Finally, the output (which was written to a file
\texttt{rptest.out}) is compared against \texttt{testbase}, which contains the
output of a sample run. The two should be identical; otherwise the
message \texttt{Test failed!} will be displayed (on its own on a line).

\begin{verbatim}
> make test1once
\end{verbatim}
runs the same curves again, but with the input fields of the argument
structure set once before the loop over the curves instead of once per
curve, as a program using the library may do since version 2.3
(Section~\ref{sec:library}); \texttt{rptest} prints a line whenever the
library has changed one of those fields, and the output is compared
against \texttt{testbase} as before.

\begin{verbatim}
> make test1many
\end{verbatim}
does the same with a second collection of curves, in
\texttt{testdata-many.h}, whose reference output is \texttt{testbase-many}.
These curves have (comparatively) many rational points, which makes them
sieve quite differently from the random curves in \texttt{testdata.h}: on a
random curve about one numerator in $10^4$ survives the first sieving stage,
on these up to a hundred times more, so the best division of labour between
the two stages is a different one. The two tests take about the same time and
between them cover the two regimes that occur in practice; both should be used
whenever the parameters are retuned, since a setting that suits one of them can
be a poor choice for the other. That is what \texttt{make tune} does with them;
see Section~\ref{sec:tuning}.

\begin{verbatim}
> make testhigh
> make testhighmany
\end{verbatim}
run the same two regimes at a height bound of $2\cdot10^5$ rather than the
$16383$ of \texttt{test1} and \texttt{test1many}, and take about a minute
each. That bound decides a good deal of what a test measures. A sieve table
is built the first time a given prime meets a given class of denominators and
is reused afterwards, so the total cost of building the tables is bounded by
the number of primes and does not grow with the height, whereas the sieving
does: on the random curves the tables are $22\%$ of \texttt{make test1} but
$0.5\%$ of \texttt{make testhigh}, and the sifting itself goes from $48\%$
to $87\%$ of the run. These are therefore the tests to judge a change to the
sieving loops by. \texttt{testhigh} uses the same curves as
\texttt{test1}; none of them has a rational point of height between $16383$
and $2\cdot10^5$, so the expected output is \texttt{testbase} again, and that
the two agree is itself worth knowing. \texttt{testhighmany} uses the thirty
curves of \texttt{testdata-many.h} that run out of sieving primes, collected
in \texttt{testdata-high-many.h} with reference output
\texttt{testbase-high-many}.

\begin{verbatim}
> make test2
\end{verbatim}
calls \texttt{ratpoints} on a curve with lots of rational points and with
a fairly large height bound, times it, and compares the output to what is
expected. This may take a few minutes, so be patient!

\begin{verbatim}
> make test3
\end{verbatim}
runs the invocations of \texttt{ratpoints} collected in \texttt{test3.sh},
one or more for each of the bugs fixed in version 2.3 (and 2.2.4), and
compares the output against \texttt{testbase3}.

\begin{verbatim}
> make timing
\end{verbatim}
times \texttt{ratpoints} on some curve with height bound $400\,000$.
This maybe useful for comparisons.

Finally,
\begin{verbatim}
> make test
\end{verbatim}
runs \texttt{test1}, \texttt{test1once}, \texttt{test1many},
\texttt{testdegrees}, \texttt{test2}, \texttt{test3} and \texttt{timing}.
The two large-height tests are not part of it; they are meant for judging a
change to the sieve rather than for checking that the build works.

\subsection{Install}

If you like, you can install the library, executable and header file
on your system.

\begin{verbatim}
> sudo make install
\end{verbatim}

The executable is copied to \texttt{/usr/local/bin/}, the
library to \texttt{/usr/local/lib/}, and the header file to
\texttt{/usr/local/include/}. You can change the \texttt{/usr/local} prefix
via
\begin{verbatim}
> make install INSTALL_DIR=...
\end{verbatim}

\subsection{Debugging}

In case you found a bug and would like to find out where it comes from,
or if you just want to see exactly what the program is doing, you can
do

\begin{verbatim}
> make debug
\end{verbatim}

in the working directory. This will build an executable \texttt{ratpoints-debug},
which, when run, will dump loads of output on the screen, so it is best
to send the output to a file, which you can then study at leisure.

\subsection{Cleaning up}

In order to get rid of the temporary files, do

\begin{verbatim}
> make clean
\end{verbatim}

To remove everything except the files from the archive, use

\begin{verbatim}
> make distclean
\end{verbatim}

\subsection{Requirements}

You need a C compiler (like \texttt{gcc}, which is the default specified
for the \texttt{CC} variable in \texttt{Makefile}; if necessary, it can be
changed there). Since version~2.3, the code assumes that \texttt{long} has
64~bits and refuses to compile otherwise; that is the case on every 64-bit
Unix-like system (but not under the Windows ABI, where \texttt{long} has
32~bits even on a 64-bit machine). Version~2.2.4 is the last one that
accommodates a 32-bit \texttt{long}.

In addition to the standard libraries, the program also requires the
\texttt{gmp} (GNU multi-precision) library~\cite{gmp}.

\subsection{List of files}

The archive contains the following files.
\begin{itemize}
  \item \texttt{Makefile}
  \item \texttt{ratpoints.h} --- the header file for programs using ratpoints
  \item \texttt{rp-private.h, primes.h} --- header files used internally
  \item \texttt{gen\_find\_points\_h.c, gen\_init\_sieve\_h.c} ---
        short programs that write additional header files depending
        on the system configuration
  \item \texttt{sift.c, init.c, sturm.c, find\_points.c} ---
        the source code for the ratpoints library
  \item \texttt{main.c, rptest.c} --- the source code for the \texttt{ratpoints}
        and \texttt{rptest} executables, respectively
  \item \texttt{testdata.h}, \texttt{testdata-many.h},
        \texttt{testdata-high-many.h}, \texttt{testdata-degrees.h} ---
        the curves used by the test runs
  \item \texttt{testbase}, \texttt{testbase-many},
        \texttt{testbase-high-many}, \texttt{testbase-degrees},
        \texttt{testbase2}, \texttt{testbase3} --- the output they are
        expected to produce
  \item \texttt{test3.sh} --- the script that \texttt{make test3} runs;
        it has to be executable
  \item \texttt{tune.sh} --- what \texttt{make tune} runs to measure the
        machine-dependent constants (Section~\ref{sec:tuning})
  \item \texttt{bench\_init.c}, \texttt{bench\_check.c} --- benchmarks for
        two parts that are timed on their own: building the sieving tables,
        and the exact test.  Neither is part of the library; the second is
        how to remeasure the constants of Section~\ref{sec:sp3choice}
  \item \texttt{ratpoints-doc-2.2.pdf} --- this documentation file
  \item \texttt{gpl-2.0.txt} --- the GNU license that applies to this program.
\end{itemize}

%=========================================================================

\section{How to use \texttt{ratpoints}}

\subsection{Basic operation}

Let
\[ C : y^2 = a_n x^n + a_{n-1} x^{n-1} + \dots + a_1 x + a_0 \]
be your
curve, and let $H$ be the bound for the denominator and absolute value
of the numerator of the $x$-coordinate of the points you want to find.
The command to search for the points is then

\verb+> ratpoints '+$a_0$ $a_1$ \dots $a_{n-1}$ $a_n$\verb+' +$H$

The first argument to \texttt{ratpoints} is the list of the coefficients,
which have to be integers, are separated by spaces, and are listed
starting with the constant term. There is no bound on the size of the
coefficients; they are read as multi-precision integers.

The degree $n$ of the polynomial is limited to \texttt{RATPOINTS\_MAX\_DEGREE}
(defined in \texttt{ratpoints.h}), which is set to~$100$ by default.

The program will give an error message when the polynomial (considered
as a binary form of the smallest even degree $\ge n$) is not squarefree
(e.g., if you specify two leading zero coefficients).

The following subsections discuss how to modify the standard behavior.
This is achieved by adding options after the two required arguments.
Some of the options have arguments, others don't; the ordering of
the options and option-argument pairs is arbitrary (except for options
with opposite effect, where the last one given counts, and for the
specification of the search intervals, which must be in order).

\subsection{Early abort}

By default, \texttt{ratpoints} will search the whole region that you specify
and print all the points it finds. If you just want to find {\em one}
point (for example when dealing with 2-covering curves of elliptic curves),
you can tell the program to quit after it has found one point via the
\verb+-1+ option, e.g.,

\verb+> ratpoints '-18 116 48 -12 30' 60000000 -1+

If you don't want to see points at infinity, specify the \verb+-i+ option.
The two options can be combined: \verb+-i -1+ will stop the program after
it has found one finite point.

\subsection{Changing the output}

By default, \texttt{ratpoints} prints some general information before and
after the points, which are given in the form
\verb+(+$x$\verb+ : +$y$\verb+ : +$z$\verb+)+, one per line.
Here, $(x : y : z)$ are the coordinates of the point considered as a point
in the $(1, \lceil n/2\rceil, 1)$-weighted projective plane, which is the
natural ambient space for the curve. The coordinates are integers with
$x$ and~$z$ coprime and $z > 0$, or $z = 0$ and $x = 1$ (for points at
infinity).

There are various ways to change this behavior.
\begin{itemize}\setlength{\parindent}{0mm}\setlength{\parskip}{1ex}
  \item To suppress all output except the points, use \verb+-q+ (for
        {\em quiet}).
  \item To add some more output explaining what the program is doing,
        use \verb+-v+ (for {\em verbose}). This has no effect if
        the \verb+-q+ option is also given.
  \item To suppress printing of the points, use \verb+-z+.
  \item If you only want to list the $x$-coordinates of point pairs rather
        than individual points, use~\verb+-y+.
  \item There are four options that influence the format in which the
        points are printed: \\
        \strut\quad \verb+-f+ {\em format} \verb+-fs+ {\em string-before}
        \verb+-fm+ {\em string-between} \verb+-fe+ {\em string-after} \\
        The arguments to all of them are strings; \verb+"\n"+, \verb+"\t"+,
        \verb+"\\"+ and \verb+"\%"+ are recognized and do what you expect.
        The {\em format} string can contain markers \verb+%x+, \verb+%y+,
        \verb+%z+ that will be replaced with the $x$, $y$, and $z$-coordinate
        of the point, respectively. The defaults are empty for
        {\em string-before}, {\em string-between} and {\em string-after},
        and \verb+"(%x : %y : %z)\n"+ for {\em format} when \verb+-y+
        is not specified; otherwise \verb+"(%x : %z)\n"+.

        The effect is as follows. Before any point is printed,
        {\em string-before} is output. Then every point is printed
        according to {\em format}. Between any two points, {\em string-between}
        is output, and after the last point, {\em string-after} is output.

        As an example, consider the effect of \\
        \strut\quad\verb+-f "[%x,%y,%z]" -fs "{" -fm "," -fe "}\n"+
\end{itemize}

\subsection{Restricting the search domain}

There are two ways to restrict the domain of the search. The first is
to restrict the range of denominators considered. This is done via \\
\strut\quad\verb+-dl+ $d_{\text{min}}$ \verb+-du+ $d_{\text{max}}$\,. \\
For example, to only look for integral points, you can say \verb+-du 1+.
By default, the lower limit is~$1$ and the upper limit is~$H$.

The other way is to restrict the search to a union of (closed) intervals.
If you only want to search for points in
\[ [l_1, u_1] \cup [l_2, u_2] \cup \dots \cup [l_k, u_k] \,, \]
you can specify this in the form \\
\strut\quad\verb+-l+ $l_1$ \verb+-u+ $u_1$ \verb+-l+ $l_2$ \verb+-u+ $u_2$ \dots
\verb+-l+ $l_k$ \verb+-u+ $u_k$ \,. \\
The first \verb+-l+ option is optional; $l_1$ defaults to~$-\infty$.
Similarly the last \verb+-u+ option is optional, and $u_k$ defaults
to~$+\infty$. The number $k$ of intervals is bounded by
\texttt{RATPOINTS\_MAX\_DEGREE} (usually, $100$).

\subsection{Setting the parameters for the sieve} \label{fd}

It is possible to change the number of primes that are used in the various
stages of the algorithm. This is done by the following options. \\
\strut\quad\verb+-p+ $M$ \verb+-N+ $N$ \verb+-n+ $n$ \\
This sets the number of (odd) primes that are considered for the sieve
to~$M$, the number of primes that are actually used for the sieve to~$N$,
and the number of primes used in the first sieving stage to~$n$. Giving~$M$
also fixes it: without \verb+-p+ the program considers
\texttt{RATPOINTS\_DEFAULT\_NUM\_PRIMES} of them to begin with and looks
further only for a curve that needs it (Section~\ref{sec:spchoice}), whereas
with \verb+-p+ it looks at exactly~$M$. The program
will, if necessary, reduce $M$, $N$, and $n$ (in this order) to ensure
that $n \le N \le M \le \text{\texttt{RATPOINTS\_NUM\_PRIMES}}$. The latter
is usually~$30$ (the number of odd primes $< 2^7$), but will be $53$
(the number of odd primes $< 2^8$) if \texttt{PRIME\_SIZE} is set to~$8$, etc.

$M$ has a price in memory, because the sieving tables for a prime~$p$ take
$p$ rows of $p$ bits and so grow with~$p^2$: the primes below $128$ need about
5\,MB, those below $1024$ about 1.7\,GB (twice that with 512-bit registers).
Only the primes that $M$ actually asks for are reserved, so raising
\texttt{PRIME\_SIZE} in order to have larger primes available costs nothing as
long as $M$ is left alone.

If \verb+-n+ or \verb+-N+ is left out, the corresponding number is chosen
from the curve, as described in Section~\ref{sec:spchoice}. The two constants
that govern that choice can themselves be set: \\
\strut\quad\verb+-r+ $x$ \verb+-R+ $e$ \\
sets the target number of numerators surviving the first sieving stage per
64-bit word to~$x$ (a real number; this is what decides~$n$), and the number
of primes added to~$n$ to obtain~$N$ to~$e$. The mnemonic is that \verb+-r+ and
\verb+-R+ are to the automatic choice what \verb+-n+ and \verb+-N+ are to
the explicit one. Both are machine-dependent, and giving them on the command
line means they can be retuned without recompiling; the compiled-in defaults
are \texttt{RATPOINTS\_SURVIVORS\_PER\_WORD}
and \texttt{RATPOINTS\_SP2\_EXTRA} in \texttt{ratpoints.h}.

Three further options control the same choice, and all three exist so that
what they switch on can be switched off again and measured: \\
\strut\quad\verb+-U+ $u$ \quad\verb+-C+ $c$ \quad\verb+-A+ $a$ \\
The first sets the length of run, in 64-bit words of numerators, at which a
second-stage prime just pays for setting itself up; \verb+-R+ is then the
number of extra primes for a run much longer than that, and a shorter run
gets proportionately fewer (Section~\ref{sec:spchoice}). \verb+-U 0+ makes
\verb+-R+ a flat offset again, as it was before version~2.3. The second sets
what building one row of a sieving table costs, in units of what one
first-stage prime costs per word, which is what lets the program prefer a
smaller prime to a larger one of the same quality; \verb+-C 0+ ranks the
primes by quality alone, as before. The third says how much of the choice the program should revise while it
runs (Section~\ref{sec:adapt}): $a = 0$ nothing, $a = 1$, the default, the
number of primes in the third stage, and $a = 2$ the number in the second
stage as well. It revises nothing it did not choose in the first place, so
\verb+-n+ or \verb+-N+ has the same effect as \verb+-A 0+.

A third sieving stage runs after those two, on the numerators that survive
them. It tests each of them against further primes one at a time, computing
the position in the table of admissible residues rather than reading a bit
out of a sieving table, so it needs no table and can use primes for which
building one would be out of the question. How many primes it uses is
chosen from the curve, and \\
\strut\quad\verb+-P+ $k$ \\
overrides that with exactly $k$ of them; \verb+-P 0+ switches the stage off,
and \\
\strut\quad\verb+-Q+ $q$ \\
sets instead the machine constant the choice turns on, namely what carrying
one more prime through the denominators costs as a fraction of one exact
test. It is to \verb+-P+ what \verb+-r+ is to \verb+-n+.
The choice weighs what a prime removes, which is a fraction of the survivors,
against what it costs, which is a little work per survivor and a little more
per denominator, so it depends on how many survivors a denominator has left
by then --- many for a curve with many rational points, few for a random one,
and more at a large height bound than at a small one. On a random curve at
the default height bound the stage does nothing at all, which is as it should
be. Section~\ref{sec:sp3choice} has the rule; the two machine constants in it
are \texttt{RATPOINTS\_SP3\_PER\_SURVIVOR} and
\texttt{RATPOINTS\_SP3\_PER\_DENOM} in \texttt{ratpoints.h}.

Both of those constants are fractions of what one exact test costs, and that
is not the same number for every curve: the test evaluates a binary form of
the given degree at one pair of integers and takes an integer square root, so
it costs more at a high degree, and more again when the coefficients are
large enough to push the arithmetic into another limb. The program estimates
it from the degree, the largest coefficient and the height bound, all of
which it knows before it looks at a single prime, and \\
\strut\quad\verb+-W+ $w$ \\
overrides the estimate with $w$ rdtsc cycles. \verb+-W 306+ is what versions
before~2.3 did in effect, namely assume that every curve costs what a
degree-6 curve with small coefficients costs; it is the way to measure what
the estimate is worth. Section~\ref{sec:sp3choice} has the formula, and
\verb+-v+ prints what it gives for the curve in hand.

There are two more parameters that can be set. \\
\strut\quad\verb+-F+ $D$ \\
sets the maximal number of `forbidden divisors'. If the degree is even,
the denominator of any $x$-coordinate of a rational point is restricted at
certain primes~$p$: it cannot be divisible by~$p$ when the leading
coefficient is a non-square mod~$p$, and when $p$ divides the leading
coefficient its valuation at~$p$ is restricted instead --- the program
works out which valuations can occur at all (Section~\ref{sec:elim}), which
typically rules out a power of~$p$, or one particular valuation, or~$p$
itself. Each such prime counts as one forbidden divisor. The program
constructs a list of them against which to check the denominators, taking
the primes up to the square root of the height bound, and this option caps
the length of the list; the default is~$64$, which the list reaches at a
height bound of about half a million. A prime in the list costs a few
instructions per 64 denominators. With \verb+-v+ it prints what it found.

The other option is \\
\strut\quad\verb+-S+ $S$ \quad or\quad \verb+-s+ \,. \\
In the first form, it sets the number of refinement (interval halving)
steps in the isolation of the real components to~$S$. This part of the
program computes a Sturm sequence for the polynomial and uses it in order
to find a union of intervals that contains the intervals of positivity
of the polynomial. This is done by a successive subdivision of the
interval $]{-\infty}, +\infty[$. The number~$S$ sets the recursion depth.
$S$ can be omitted; then it is given a default value. The option
\verb+-s+ skips the Sturm sequence computation completely. This also has
the effect of removing most of the check for squarefreeness.

\subsection{Switching off optimizations}

The options \verb+-k+ and \verb+-j+ can be used to prevent the program
from reversing the polynomial, which it usually does if this will lead
to faster operation (\verb+-k+), and to prevent the program from using
the Jacobi symbol test on the denominators (\verb+-j+). This last test
extends the `forbidden divisors' method described in
Subsection~\ref{fd} above by computing the Jacobi symbol $(l/d)$, where
$l$ is the leading coefficient and $d$ is the odd and coprime-to-$l$
part of the denominator. If the symbol is~$-1$, the denominator need not
be considered. The use of these options may be questionable, unless you
want to see how much performance is gained by using these optimizations.

\subsection{Switching off exact testing of points}

The option \verb+-x+ will prevent the program from checking the potential
points that survive the sieve whether they really give rise to rational
points. This implies that the points that are output may not actually
be points. It also effectively sets the \verb+-y+ option, since the
$y$-coordinates are not computed.

\subsection{Overriding previous options}

The options \verb+-I+, \verb+-Y+, \verb+-Z+, \verb+-S+, \verb+-K+, \verb+-J+,
\verb+-X+ can be used to cancel the effect of the corresponding lower-case
option (and thereby restore the default behavior),
when this occurs earlier in the list of options. This may be useful if
you want to set a default behavior that is different from what the
program does out-of-the-box (e.g., in a shell script), but want the caller
to be able to override this change.

%=========================================================================

\section{How to use the library}\label{sec:library}

It is possible to use the ratpoints machinery from within your own programs.

The library \texttt{libratpoints.a} provides the following functions.

\begin{verbatim}
long find_points(ratpoints_args*,
                 int proc(long, long, const mpz_t, void*, int*),
                 void*);

void find_points_init(ratpoints_args*);

long find_points_work(ratpoints_args*,
                      int proc(long, long, const mpz_t, void*, int*),
                      void*);

void find_points_clear(ratpoints_args*);
\end{verbatim}

The passing of arguments to these functions is via the \texttt{ratpoints\_args}
structure, which is defined as follows.

\begin{verbatim}
typedef struct { mpz_t *cof; long degree; long height;
                 ratpoints_interval *domain; long num_inter;
                 long b_low; long b_high; long sp1; long sp2; long sp3;
                 double survivors_per_word; long sp2_extra; double sp2_u0;
                 double cost_table; long adapt;
                 long sp3_extra; double sp3_per_denom; double check_cost;
                 long array_size;
                 long sturm; long num_primes; long max_forbidden;
                 unsigned int flags;
                 long sp1_used; long sp2_used; long sp3_used; ...}
        ratpoints_args;
\end{verbatim}

The dots at the end stand for additional fields that are used internally
and are not of interest here. All fields up to and including \texttt{flags}
are inputs, and they come back from a call as they went in, with three
exceptions that are described below: the coefficient array and
\texttt{degree}, which on return describe the polynomial the program
worked with; \texttt{num\_inter} and \texttt{domain}, which on return
describe the region it searched; and the flag bits that report on the
run. (Up to version 2.2.4 the program also wrote the values it chose or
normalised into the other input fields, so that a caller who fills the
structure once and then loops over curves had to reset them before every
call; this is no longer necessary.) Every input field has to be set before
the call; a negative value asks for the compiled-in default, or for the
value to be chosen from the curve, as described below.
The three \texttt{\_used} fields are outputs, see below.

When calling \texttt{find\_points}, the
\texttt{cof} field must point to an array of size $\text{\texttt{degree}} + 1$
of properly initialized gmp integers; these are the coefficients of
the polynomial. The program can alter the values of the integers in this
array: it may reverse the polynomial, and it drops a zero constant term
on that occasion, in which case \texttt{degree} is lowered by one so that
\texttt{cof} and \texttt{degree} still describe the polynomial the program
worked with (the \texttt{RATPOINTS\_REVERSED} flag says when this has
happened). Leading zero coefficients are dropped in the same way, lowering
\texttt{degree}, without a flag.
To prevent the reversal, set the \texttt{RATPOINTS\_NO\_REVERSE} bit
in \texttt{flags}; this may result in a loss of performance, though.

\texttt{height} gives the height bound. It is an error for the
degree or the height bound to be nonpositive; in this case the function
returns the value \texttt{RATPOINTS\_BAD\_ARGS}.

The field \texttt{domain} must contain a pointer to an array of
\texttt{ratpoints\_interval} structures of length at least \texttt{num\_inter}
plus \texttt{degree}. This array gives (in its first \texttt{num\_inter} entries)
the intervals for the search region. The type \texttt{ratpoints\_interval}
is just

\begin{verbatim}
typedef struct {double low; double up;} ratpoints_interval;
\end{verbatim}

the meaning of this should be clear. In \texttt{ratpoints}, this is set
via the \verb+-l+ and \verb+-u+ options. Usually, you do not want to
restrict the range of $x$-coordinates; then you set \texttt{num\_inter = 0}
(but you still have to fill \texttt{domain} with a valid pointer to at least
\texttt{degree} intervals; a null pointer makes the call return
\texttt{RATPOINTS\_BAD\_ARGS}, since the program stores the default
interval $]-\texttt{height}, \texttt{height}[$ there. If \texttt{sturm} is
negative, one interval is enough.)
The program intersects the given intervals with the region where $f$ is
positive (unless \texttt{sturm} has a negative value, which may lead to a
performance loss) and searches the result; on return, \texttt{num\_inter}
and the first \texttt{num\_inter} entries of the array describe that
result, which is why the array has to be longer than the input. So
\texttt{num\_inter} (and \texttt{domain}, if you restrict the search) must
be set before every call.

\texttt{b\_low} and \texttt{b\_high}
carry the lower and upper bounds for the denominator; if non-positive, they are
set to~$1$ and the height bound, respectively. In \texttt{ratpoints}, these
fields are set by the \verb+-dl+ and \verb+-du+ options.

\texttt{sp1} and \texttt{sp2}
specify the number of primes to be used in the first sieving stage and
in both sieving stages together, respectively. If negative, they are chosen
from the curve, as described in Section~\ref{sec:spchoice}; this is what the
command line program does when \verb+-n+ and \verb+-N+ are absent. In
\texttt{ratpoints}, these fields are set by the
\verb+-n+ and \verb+-N+ options. What was actually used, chosen or not,
is reported in \texttt{sp1\_used} and \texttt{sp2\_used}. The fields
\texttt{survivors\_per\_word} and \texttt{sp2\_extra} hold the two
constants that govern that choice (options \verb+-r+ and \verb+-R+); if
negative, the values compiled into \texttt{ratpoints.h} are used. So do
\texttt{sp2\_u0} and \texttt{cost\_table} (options \verb+-U+ and \verb+-C+),
which say how long a run has to be for a second-stage prime to pay for
setting itself up and what one row of a sieving table costs; zero in either
switches off the corresponding refinement. \texttt{adapt} set to zero stops
the program revising these numbers as it runs (option \verb+-A+,
Section~\ref{sec:adapt}); it revises them only when it chose them, so a
non-negative \texttt{sp1} or \texttt{sp2} has the same effect.

\texttt{sp3\_used}
reports the number of primes used in all three stages together, so that
\texttt{sp3\_used} $-$ \texttt{sp2\_used} of them belong to the third.
(\texttt{sp3} itself is a working field of no interest to the caller.)
What goes in is \texttt{sp3\_extra}, which fixes that difference when
it is non-negative (option \verb+-P+) and otherwise leaves the number to be
chosen from the curve, and \texttt{sp3\_per\_denom} (option \verb+-Q+),
which is the machine constant that choice turns on and takes the value
compiled into \texttt{ratpoints.h} when negative. \texttt{check\_cost}
(option \verb+-W+) says what one exact test costs, in rdtsc cycles, and both
of those constants are fractions of it; when negative, which is the usual
case, the program estimates it from the degree, the coefficients and the
height bound as described in Section~\ref{sec:sp3choice}.
Similar statements are true for
\texttt{num\_primes} (option \verb+-p+), \texttt{sturm} (options \verb+-s+,
\verb+-S+) and \texttt{max\_forbidden} (option \verb+-F+). The field
\texttt{array\_size} specifies the maximal size (in bit arrays; a bit array
contains 64, 128, 256 or 512 bits, depending on the configuration) of the
array that is used in the first sieving stage. If non-positive, it is
set to a default value. The various default values are defined at the
beginning of the header file \texttt{ratpoints.h}, where also the maximal
degree of the polynomial is set.

The \texttt{flags} field holds a number of bit flags.
\begin{itemize}
  \item \texttt{RATPOINTS\_NO\_CHECK} --- when set, do not check whether
        the surviving $x$-coordinates give rise to rational points
        (set by the \verb+-x+ option to \texttt{ratpoints}).
  \item \texttt{RATPOINTS\_NO\_Y} --- only list $x$-coordinates (in the
        form $(x : z) \in \BP^1(\Q)$) instead of actual points (with
        a $y$-coordinate); this is set by the \verb+-y+ option to
        \texttt{ratpoints}.
  \item \texttt{RATPOINTS\_NO\_REVERSE} --- when set, do not allow reversal
        of the polynomial (set by the \verb+-k+ option to \texttt{ratpoints}).
  \item \texttt{RATPOINTS\_NO\_JACOBI} --- when set, prevent the use of the
        Jacobi symbol test (set by the \verb+-j+ option to \texttt{ratpoints}).
  \item \texttt{RATPOINTS\_VERBOSE} --- when set, causes the procedure to
        print some output on what it is doing (set by the \verb+-v+ option
        to \texttt{ratpoints}).
\end{itemize}

There are some other flags that are used internally. One of them might be
of interest:
\begin{itemize}
  \item \texttt{RATPOINTS\_REVERSED} --- when set after the function call,
        this indicates that the polynomial has been reversed (and the
        contents of the \texttt{cof} array have been modified).
\end{itemize}

The main vehicle for passing information back to the caller is the
\texttt{proc} function argument together with the pointer \texttt{info}.
This function \\
\strut\quad\verb+int proc(long x, long z, const mpz_t y, void *info, int *quit)+ \\
is called whenever a point was found. \texttt{x}, \texttt{y} and \texttt{z} are
the coordinates of the point (where \texttt{y} is a gmp integer).
\texttt{info} is the pointer that was passed to \texttt{find\_points}; this
can be used to store information that should persist between calls
to \texttt{proc}. If \texttt{*quit} is set to a non-zero value, this indicates
that \texttt{find\_points} should abort the point search and return immediately;
otherwise the search continues. This is how the \verb+-1+ option works.
The return value is taken as a weight
for counting the points; usually it will be~1.

The usual framework for using \texttt{find\_points} is as follows.

\begin{verbatim}
(...)
#include "ratpoints.h"

(...)

mpz_t c[RATPOINTS_MAX_DEGREE+1];  /* The coefficients of f */
ratpoints_interval domain[2*RATPOINTS_MAX_DEGREE];
                     /* This contains the intervals representing the
                        search region */

/*****************************************************************
 * function that processes the points                            *
 *****************************************************************/

typedef struct {...} data;

int process(long x, long z, const mpz_t y, void *info0, int *quit)
{ data *info = (data *)info0;

  (...)

  return(1);
}
\end{verbatim}


\begin{verbatim}
/*****************************************************************
 * main                                                          *
 *****************************************************************/

int main(int argc, char *argv[])
{
  long total, n;
  ratpoints_args args;

  long degree        = 6;
  long height        = 16383;
  long sieve_primes1 = RATPOINTS_DEFAULT_SP1;
  long sieve_primes2 = RATPOINTS_DEFAULT_SP2;
  long num_primes    = RATPOINTS_DEFAULT_NUM_PRIMES;
  long max_forbidden = RATPOINTS_DEFAULT_MAX_FORBIDDEN;
  long b_low         = 1;
  long b_high        = height;
  long sturm_iter    = RATPOINTS_DEFAULT_STURM;
  long array_size    = RATPOINTS_ARRAY_SIZE;
  int no_check       = 0;
  int no_y           = 0;
  int no_reverse     = 0;
  int no_jacobi      = 0;
  int no_output      = 0;

  unsigned int flags = 0;

  data *info = malloc(sizeof(data));

  /* initialize multi-precision integer variables */
  for(n = 0; n <= degree; n++) { mpz_init(c[n]); }

  (...)

  { /* set up polynomial */
    long k;
    for(k = 0; k < 7; k++) { mpz_set_si(c[k], ...); }

    args.cof           = &c[0];
    args.degree        = 6;
    args.height        = height;
    args.domain        = &domain[0];
    args.num_inter     = 0;
    args.b_low         = b_low;
    args.b_high        = b_high;
    args.sp1           = sieve_primes1;
    args.sp2           = sieve_primes2;
    args.survivors_per_word = -1.0; /* negative: the compiled-in default */
    args.sp2_extra     = -1;
    args.sp2_u0        = -1.0;
    args.cost_table    = -1.0;
    args.adapt         = -1;
    args.sp3_extra     = -1;        /* negative: chosen from the curve */
    args.sp3_per_denom = -1.0;
    args.check_cost    = -1.0;
    args.array_size    = array_size;
    args.sturm         = sturm_iter;
    args.num_primes    = num_primes;
    args.max_forbidden = max_forbidden;
    args.flags         = flags;

    info->... = ...;
    (...)

    total = find_points(&args, process, (void *)info);
    if(total == RATPOINTS_NON_SQUAREFREE)
    { ... }
    if(total == RATPOINTS_BAD_ARGS)
    { ... }

    (...)

  }

  /* clean up multi-precision integer variables */
  for(n = 0; n <= degree; n++) {mpz_clear(c[n]); }

  return(0);
}
\end{verbatim}

If points are to be searched on many curves of the same degree,
then it is slightly more efficient to use the sequence

\begin{verbatim}
  args.degree = degree; /* this information is needed */
  find_points_init(&args);

  for( ... )
  { ...
    total = find_points_work(&args, process, (void *)info);
    ...
  }

  find_points_clear(&args);
\end{verbatim}

This avoids the repeated allocation and freeing of memory. Inside the loop
only \texttt{cof}, \texttt{degree} and \texttt{num\_inter} (and whatever
else changes from curve to curve) need to be set; the other input fields
keep their values across calls, since version 2.3. \texttt{rptest.c} with
the \verb+-O+ option does exactly this, and \texttt{make test1once} checks
that the points found are the same as with the fields set once per curve.

For practical examples, see \texttt{main.c} (the code that wraps
\texttt{find\_points} for the command line program \texttt{ratpoints}) or
\texttt{rptest.c} (which runs \texttt{find\_points} on some test data).

%=========================================================================

\section{Fine-tuning the parameters}\label{sec:tuning}

\subsection{How \texttt{sp1} and \texttt{sp2} are chosen}\label{sec:spchoice}

Unless they are given explicitly, the two most important parameters are
chosen from the curve. While collecting the sieving information, the program
computes, for each prime~$p$, the proportion~$r_p$ of residues~$a$ modulo~$p$
for which $f(a)$ is a square, and sorts the primes by increasing~$r_p$; the
expected proportion of numerators surviving the first $k$~of them is then
the product of the $k$~smallest~$r_p$. Multiplied by the number of bits that
are set in a bit array to begin with, this is the expected number of
survivors per bit array.

The first sieving stage costs one load and one \texttt{AND} per bit array
per prime, whatever survives, while the second stage costs nothing on an
empty bit array and a good deal on a non-empty one. So primes should be moved
into the first stage until few enough bit arrays are left alive. Measurements
over curves spanning a factor of~$600$ in density show that the best
\texttt{sp1} is, to within a factor of about two, the one at which the
expected number of survivors per 64-bit word falls below a constant --- almost
independently of the curve, even though the resulting \texttt{sp1} itself
ranges from $8$ to~$30$. That constant is
\texttt{RATPOINTS\_SURVIVORS\_PER\_WORD} in \texttt{ratpoints.h}; it is a
property of the machine rather than of the curve. Per word rather than per bit
array on purpose: tuning runs at 128, 256 and 512 bits give $0.0075$, $0.0075$
and $0.010$ survivors per word, so the same constant serves every register
width, whereas a per-bit-array figure would have to double with each doubling
of the width. Then \texttt{sp2} is set to
\texttt{sp1} plus \texttt{RATPOINTS\_SP2\_EXTRA}, capped by the number of
primes that carry any information at all.

That offset is not the same for every run, because part of what a prime costs
is paid once and then spread over the whole run. Its sieving table is built
at most $p$ times, however many numerators there are, and its entry in the
list of denominators reduced modulo the sieving primes is computed once per
denominator; per word of numerators both fall like $1/U$, where $U$ is the
number of numerator words the run will sweep. The program computes $U$, and
the number $D$ of denominators, before it looks at a single prime --- the
numerators of one denominator are a piecewise linear function of it, so a
sample over the range of denominators is accurate to a fraction of a per cent
--- and takes
\[ \texttt{sp2} = \texttt{sp1}
   + \frac{\texttt{RATPOINTS\_SP2\_EXTRA}}{1 + \texttt{RATPOINTS\_SP2\_U0}/U} \,. \]
\texttt{RATPOINTS\_SP2\_U0} is the length of run at which setting a prime up
costs as much as sieving with it; like the other two it is a property of the
machine. Setting it to zero restores a flat offset. This is what lets one
tuning serve every height bound. Sweeping the offset gives four measured
optima --- $3$ on random curves at a height bound of $16383$, $5$ on
point-rich ones there, $9$ on random curves at $2\cdot 10^5$ and $12$ on
point-rich ones --- and the two constants land on $3$, $8$, $11$ and $11$.
Using the small-bound value at the large one costs $12\%$ of \texttt{make
testhighmany}. The first stage needs no such correction, and measurement
agrees that it does not: a first-stage prime costs one \texttt{AND} per word
unconditionally, so the fixed part is a far smaller share of what it costs,
and the best threshold hardly moves with the height bound.

The one thing this shape cannot fit is the point-rich curves at the small
height bound, which want $5$ and are given $8$. Their survivor rate is twice
the random curves', so a second-stage prime costs them twice as much per
word, and a formula in the length of the run alone cannot see that.

The same two quantities say which primes to use, not just how many. The
sieving table for~$p$ has $p$ rows and is rebuilt for each denominator class
that turns up, at most $p$ of them, so it costs $O(p^2/U)$ per word: at equal
quality a smaller prime is strictly better, and at a small height bound it is
much better. The primes are therefore ranked not by~$r_p$ but by what they
cost divided by what they say,
\[ \frac{\text{cost}(p)}{-\log r_p} \,, \]
and ranked twice, once for each of the first two stages, because a
second-stage prime is applied only to the bit arrays that survived the first
and its table therefore weighs some three times more heavily against it. The
third stage builds no table, so there the ranking is by $r_p$ alone, as it
always was. \texttt{RATPOINTS\_COST\_TABLE} is what one row of a table costs,
in units of what one first-stage prime costs per word; setting it to zero
makes the ranking monotone in~$r_p$ and so restores the old order exactly.

That cap is a real constraint on a curve with very many rational points, and
the program does something about it. On such a curve $f$ is a square modulo
every residue for the smallest primes, so those primes say nothing and are
discarded, and there may be fewer than $\texttt{sp1} + \texttt{RATPOINTS\_SP2%
\_EXTRA}$ left --- in the extreme, exactly \texttt{sp1} of them, so that the
second sieving stage has no primes to sieve with and every survivor of the
first stage goes to the exact test. So when the primes it was given do not
suffice, the program looks at further primes, one at a time, until they do or
until there are no more. Adding a prime can only lower \texttt{sp1}, since the
$k$ smallest of a larger set have a smaller product, and can only raise the
number of informative primes, so the shortfall shrinks with every prime added
and the search stops at the first prime that is enough. It is asked for only
when it is needed: of the $98$ curves in \texttt{testdata-many.h} it leaves
$68$ untouched and takes at most six extra primes for the rest, and of the
random curves in \texttt{testdata.h} not one triggers it. On the curves that
do need it, it is worth between $1.3$ and $2$ times at a height bound
of~$2\cdot 10^5$. Giving \verb+-p+ explicitly switches this off: an explicit
number of primes is a hard limit.

Because \texttt{sp1} is chosen this way, the second sieving stage now finds
only a few per cent of the bit arrays non-empty whatever the curve, and it
spends most of its time stepping over the empty ones. It does that in a tight
loop of its own rather than re-entering the main loop body each time, which is
worth some $7\%$ on its own. Making that loop test several
bit arrays at once, by or-ing them together, used to pay when \texttt{sp1} was
fixed and the proportion of surviving bit arrays therefore varied a great deal;
now that the proportion is held down to a couple of per cent it no longer does,
since a group containing a survivor wastes its whole `or', and that happens
often enough to cost more than the tests it saves.

Taken together, and against the fixed values used up to version~2.2.3, this is
about $14\%$ faster on the random curves of \texttt{make test1} and about
$36\%$ faster on the point-rich curves of \texttt{make test1many}.

\subsection{How many primes the third stage uses}\label{sec:sp3choice}

The third stage tests one survivor at a time, so what a prime costs it and
what a prime saves it are of quite a different shape from the first two
stages. Saved is a fraction $1 - r_p$ of the survivors that reach the prime,
each of which would otherwise go to the exact test with its multi-precision
arithmetic and its integer square root. Spent is a multiplication and a
reduction for each survivor tested, and one step of the loop that keeps the
denominator modulo each sieving prime --- and that last is paid for every
denominator, whether or not it has a survivor left. So a prime is worth
adding as long as
\[ S\,(1 - q_s - r_p) > q_d , \]
where $S$ is the expected number of survivors a denominator still has and
$q_s$, $q_d$ are the per-survivor and per-denominator costs measured against
the cost of one exact test. Since $S$ falls by a factor of~$r_p$ with each
prime added, this stops of its own accord. The two constants are
\texttt{RATPOINTS\_SP3\_PER\_SURVIVOR} and
\texttt{RATPOINTS\_SP3\_PER\_DENOM} in \texttt{ratpoints.h}; like the two
above they are properties of the machine.

They are fractions of one exact test, and what one exact test costs is a
property of the curve rather than of the machine alone. The test computes
\[ F(a,b) = c_d a^d + c_{d-1} a^{d-1} b + \dots + c_0 b^d \]
by Horner's rule from the $c_k b^{d-k}$, which are formed once per
denominator, and then takes the integer square root of the result. So it
costs one multiplication by a single word and one addition per step of the
Horner loop, that is one per degree, on numbers that grow from one limb to
the size of $F$; one more multiplication when the degree is odd, to make the
form of even degree; and a square root, which is nearly free while the root
fits in one limb and costs a fixed amount for each further limb of the root
after that. Since $F(a,b)$ has about $\log_2|c| + d\log_2 H$ bits, the degree,
the largest coefficient and the height bound fix every quantity in that
count, and the program forms the estimate before it looks at a prime. Both
fractions are then divided by the estimate over its value for a degree-6
curve with small coefficients, so they are unchanged on such a curve, smaller
on one whose test costs more, and larger on one whose test costs less --- and
a curve with an expensive test is worth more third-stage primes, because each
one it adds is spared more.

The four constants of the estimate are \texttt{RATPOINTS\_CHECK\_STEP},
\texttt{RATPOINTS\_CHECK\_LIMB}, \texttt{RATPOINTS\_CHECK\_CALL} and
\texttt{RATPOINTS\_CHECK\_ROOT}, with the reference value in
\texttt{RATPOINTS\_CHECK\_REFERENCE}; only the ratio of one curve's test to
another's is used, so a machine on which every test is dearer needs no change
to any of them. On this machine the estimate agrees with what the sieve
itself shows to within about six per cent over degrees 3 to~12 and
coefficients of 4 to 120 bits, the exceptions being the curves whose $F$ sits
within a few bits of a limb boundary, where the square root may or may not
take the further limb.

$S$ is estimated from the number of numerators a denominator has to consider
--- which follows from the positivity intervals, the height bound and the
range of denominators --- thinned by the sixteen-fold pre-sieve, by the
primes of the first two stages, and by the proportion of numerators that are
coprime to the denominator. That last factor is there because the test for
common factors runs before this stage and not after it. It has to: a
numerator sharing a factor with the denominator stands for a fraction that a
smaller denominator has already dealt with, so $f$ takes the same value there
and every prime accepts it. Those survivors are shadows of the curve's own
rational points, and there are about $H/b$ of them for each point with
denominator~$b$, so at a small height bound they can outnumber the survivors
that die by chance several times over.

The stage also decides how far up the prime list to look. A random curve has
a dozen or more informative primes left over past \texttt{sp2} among the
thirty the program looks at by default, but a curve with many rational points
may have none at all --- these are the curves of the previous section, whose
small primes say nothing. Since the stage builds no sieving table, a prime
costs it nothing per curve beyond the table of admissible residues, so it
looks at further primes when the ones in hand run out, exactly as the choice
of \texttt{sp2} does. On the curves of \texttt{testdata-high-many.h} that
turns one available prime into more than twenty.

\subsection{Correcting the choice while the program runs}\label{sec:adapt}

Everything above is decided before a single numerator has been sieved, from
the densities~$r_p$ alone, and that prediction is wrong in a way no amount of
care will put right. If $a/b$ is a rational point, so is $ka/kb$ for every~$k$
--- it is the same rational number --- and it gives the same value of~$f$,
so it passes every prime test there is. A floor of survivors therefore
outlives any amount of sieving, and only the test for common factors removes
it. The prediction cannot see that floor, so it always overshoots; that is
why \texttt{sp2} could never be predicted outright and ended up as an offset
with a constant to tune.

The program can see it, though, because it is sieving. The scan visits every
bit array anyway, so counting the non-empty ones costs an increment on a path
taken half a per cent of the time, and the survivors of the second stage and
of the test for common factors are counted as cheaply. Writing
\[ S(n) = \text{floor} + \text{chance}\cdot R(n) \]
for the survivors per word after $n$~primes, with $R(n)$ the product of the
$n$~smallest~$r_p$, two such counts pin both terms. What is done with that
comes in two degrees. By default the program corrects only the third stage,
where the number of survivors a denominator brings to the stage becomes a
count rather than an estimate, and where nothing else supplies that number.
With \verb+-A 2+ it also reconsiders \texttt{sp2} by the same marginal rule
--- which is a second answer to a question the scaled offset above already
answers, so the two are not meant to be believed at once; whichever is
applied later wins, and that is the correction. The first correction is made after a million
numerator words have been swept and the next whenever twice as many have
been, so the estimate improves and the corrections thin out.

Changing the number of primes part-way through a run is safe by construction.
Sieving only ever removes numerators that cannot be points, so using more or
fewer primes for later denominators changes how long the program takes and
nothing whatever about what it finds. This only happens when the program
chose the numbers itself: \verb+-n+, \verb+-N+ or \verb+-A 0+ leaves them
alone.

\subsection{Retuning for your machine}

For large computations, it may be a good idea to try to find the best
(or at least, a good) combination of parameters for the given kind of data.
The four constants above --- the threshold, the offset, the run length at
which a second-stage prime pays for setting itself up, and what one row of a
sieving table costs --- were tuned on my current laptop by minimising the
sum of the times of \texttt{make test1} and \texttt{make test1many}. To repeat
that on your machine, run
\begin{verbatim}
> make tune
\end{verbatim}
which does the whole thing: it sweeps the threshold, then the offset, then
the run length at which a second-stage prime pays for setting itself up, and
last what one row of a sieving table costs, over both test sets, and writes
the values it finds to a file
\texttt{tuning.mk} that the \texttt{Makefile} includes, so that a subsequent
\texttt{make all} builds with them. No source file is touched; deleting
\texttt{tuning.mk} restores the values compiled into \texttt{ratpoints.h}.
The file records the configuration it was measured for, and is ignored (with a
warning) if you later build with a different register width or
\texttt{PRIME\_SIZE}, since the best values may differ.

It is deliberately not part of \texttt{make all}: it takes several minutes,
wants an otherwise idle machine, and must not be run under \texttt{make -j}.
How long depends on the register width, since that is what decides how long the
two tests take --- on the machine this was developed on, about three minutes
for the 256-bit build, four for 128-bit, six for the plain 64-bit one and eight
for the 512-bit code emulated through narrower operations.

Four environment variables adjust it, should the defaults not suit:
\texttt{ROUNDS} (how many times each setting is measured, default~$3$),
\texttt{WARMUP} (seconds of load before anything is counted, default~$20$),
\texttt{MARGIN} (how much better the winner must be, default $0.97$) and
\texttt{NOISE} (how much residual noise is tolerated, default $0.02$). For
instance \texttt{ROUNDS=6 make tune} on a machine that reports itself too
noisy.

Which tests it measures on is worth a thought. At the height bound
\texttt{test1} and \texttt{test1many} use, a fifth of the time on the random
curves goes into building sieve tables, which the threshold and the offset
have no effect on (the table cost is the one constant that bears on it), so a
setting is judged partly on work it does not touch. \texttt{make
tunehigh} measures on the large-height tests instead, where the sieve is nearly
all of the run; use it if the runs that matter to you are long ones.

One timing there is two minutes rather than three seconds, so it does not sweep
the whole ladder of candidates. It starts from the settings in force --- which
is to say from what \texttt{make tune} found, since the two read and write the
same \texttt{tuning.mk} --- and asks only whether a factor of two in the
threshold, in the run length or in the table cost either way, or two more or
fewer primes in the second stage, is better. That is fifteen settings in a
round rather than twenty-four,
and it rests on the assumption that the two regimes do not want wildly
different values. If it moves a value it has not finished looking, and should
be run again from there. Expect a couple of hours at the default three rounds;
\texttt{ROUNDS=1 make tunehigh} is the short version. The intended order is
\texttt{make tune} first, for a value from a full sweep cheaply, and then
\texttt{make tunehigh} to ask whether the long runs want something else.

Running both is also the only way to pin the third constant. The offset is
now the number of extra primes an arbitrarily long run wants, and a given run
gets a fraction of that which depends on how long it is, so a single tuning
sees only the product and cannot separate the two. Two tunings at run lengths
a hundred times apart can.

The candidates themselves can be chosen too. \texttt{R\_VALUES},
\texttt{E\_VALUES}, \texttt{U\_VALUES} and \texttt{C\_VALUES} give an
absolute ladder, which is what \texttt{make tune} uses, and
\texttt{R\_FACTORS}, \texttt{E\_DELTAS}, \texttt{U\_FACTORS} and
\texttt{C\_FACTORS} give a neighbourhood of the current settings ---
multiples of the threshold, of the run length and of the table cost, and
offsets added to the number of primes --- which is what \texttt{make
tunehigh} uses and which takes precedence when set.

Timing this reliably is harder than it looks, and \texttt{tune.sh} goes to
some trouble over it. The cost surface is flat --- anything within a factor of
two of a good threshold costs under $3\%$ --- whereas a laptop under sustained
load slows by a quarter as it warms up, so a straight sweep would measure the
clock rather than the settings. Each candidate is therefore timed back to back
with the current settings and only the ratio is kept, the median over several
rounds is used, and the current settings are themselves among the candidates,
so that the ratio they measure against themselves shows how much noise is
left. If that is more than $2\%$ the run is declared inconclusive, and
nothing is written unless the winner is better by more than $3\%$ --- a noisy
run should leave a good default alone rather than replace it with a worse one.

Two consequences are worth expecting rather than being surprised by. A run
that reports that nothing beat the current settings has \emph{succeeded}: on a
build that is already tuned that is the right answer, and it is what happens
here. And a configuration whose tests take long enough for the machine's
thermal behaviour to change within a single round --- the plain 64-bit build on
a laptop, for instance --- may well report itself inconclusive however often it
is run; raising \texttt{ROUNDS}, or using a machine that does not throttle,
is then the remedy.

Should you prefer to do it by hand, the constants can be set directly:
\begin{verbatim}
> ./rptest -r 0.0075 -R 11 -U 1.6e6 -C 38 -z
> ./rptest-many -r 0.0075 -R 11 -U 1.6e6 -C 38 -z
\end{verbatim}
for a few values of each, minimising the sum of the two times; then put the
values you found into \texttt{ratpoints.h}. The option \verb+-T+ makes the
test print the CPU time it used, and nothing else, which is what
\texttt{tune.sh} uses. The threshold \verb+-r+ matters considerably more
than the offset \verb+-R+, which at large height bounds hardly matters at
all. Use both tests: they cover the two regimes that occur in practice, and a
value that suits one of them can be a poor choice for the other.

For your machine
and input data, other values may be better. You can replace the
\texttt{testdata.h} file with your own collection of test data (and
edit \texttt{rptest.c} if necessary to adapt the degree and height settings)
and then time \texttt{./rptest -z} (the \verb+-z+ option suppresses the output)
for various combinations of the parameter settings. \texttt{rptest} accepts
most of the optional parameters of \texttt{ratpoints}, in particular
the \verb+-p+, \verb+-N+, \verb+-n+, \verb+-r+, \verb+-R+, \verb+-U+,
\verb+-C+, \verb+-P+, \verb+-Q+, \verb+-F+ and \verb+-S+ parameters, as well as \verb+-T+ (print the CPU time used, and
nothing else),
so you can easily change the parameters used on the command line. In addition,
there is an option \verb+-h+ $H$ to change the default height bound
and an option \verb+-m+ $m$ that causes the program to repeat the
computations $m$~times.

Once you have found a good set of parameter values for your application,
you can hard-code them as defaults into \texttt{ratpoints} by changing the
definitions in \texttt{ratpoints.h} (and then \texttt{make distclean all test}),
or you can use them to fill the \texttt{ratpoints\_args} structure for
your call to \texttt{find\_points}.

%=========================================================================

\section{Implementation} \label{impl}

\subsection{Overview}

Let $F(x, z)$ be the binary form of even degree corresponding to the
polynomial on the right hand side of the curve equation. The basic idea
is to let run $b$ from $1$ to the height bound~$H$, for each~$b$, let
$a$ run from $-H$ to~$H$, and for each coprime pair $(a,b)$ check
if $F(a, b)$ is a square.

Of course, in this form, this would take a very long time. To speed up the
process, we try to eliminate quickly as many pairs $(a, b)$ as possible
before the actual test. This can be done by `quadratic sieving' modulo
several primes: if $F(a, b)$ is a square, it certainly has to be a square
mod~$p$, and so we can rule out all $(a, b)$ that do not satisfy this
condition. Another ingredient is to represent (for a fixed~$b$) the
various values of~$a$ by bits and treat all the bits in a `bit array'
(of 64 bits, or of 128, 256 or 512 bits when SSE/AVX instructions are
used) in parallel. For this, we organize the sieving information for each
prime~$p$ into $p$~arrays of $p$~words each, one such array for every $b$
mod~$p$, such that the $j$th~bit in this array is set if and only if
$F(j, b)$ is a square mod~$p$.

For each `denominator' $b$, we then set up an array of words whose bits
represent the range $-H \le a \le H$ (aligned so that $a = 0$ corresponds
to the $0$th bit of a word); the bits are initially set. Then for each
of the sieving primes~$p$, we perform a bit-wise {\em and} operation
between this array and the sieving information at~$p$. In a first stage,
this is done on the whole array; after this first stage, each remaining
(`surviving') non-zero word in the array is subject to tests with more
primes. If some bits are still set after this second stage of sieving,
the corresponding pairs $(a, b)$ (if coprime) are then checked exactly.

In the following subsections, we discuss a number of improvements that
were made.

\subsection{Sorting the primes}

This idea is due to \textbf{Colin Stahlke}. We do not just take the first so many
primes in increasing order for the sieving, but we first compute the
number of points the curve has mod~$p$ for a number of primes~$p$ and then
sort the primes according to the fraction of $x$-coordinates that give points.
We then take those primes for the sieving that have the smallest fraction
of `surviving' $x$-coordinates. In this way, we need fewer sieving primes
to achieve a comparable reduction of point tests.

\subsection{Using connected components}

We note that we can only have points when $F(a, b)$ is non-negative.
If we can determine intervals on which $f(x) = F(x, 1)$ is negative, then
we do not have to look for points in these intervals. The necessary
computations can be performed exactly, by computing a Sturm sequence
for~$f$ and counting the number of sign changes at various points,
see~\cite[Thm.~4.1.10]{Cohen}. This can in particular tell us whether
$F$ is negative definite, in which case we have already proved that there
are no rational points in the curve. If there are real zeros, we use
a subdivision method in order to find a collection of intervals containing
the projections of the connected components of the curve over~$\R$.

\subsection{Using 2-adic information}

\textbf{John Cremona} suggested that in some cases, one can determine beforehand
that all points will have odd `numerators'~$a$, and so we can pack the
bits more tightly by only representing odd numbers. This approach can be
extended. We first find all solutions mod~$16$ (higher powers of~$2$ would
be possible, but not very likely to give a significant improvement). Then for every
residue class mod~$16$ of the denominator~$b$, we can find the residue
classes mod~$16$ of potential numerators~$a$. If there are none, then we
can eliminate $b$ as a denominator altogether. If all potential $a$'s are
even or odd (this will always be the case for even denominators), we can
restrict the sieving to such numerators.

\subsection{Elimination of denominators}

Depending on the equation, certain denominators can be excluded. We have
seen an instance of this in the previous subsection, but we can also work
with odd primes.

\subsubsection{Odd degree and $\pm$monic}

If the polynomial has odd degree and leading coefficient $\pm 1$, then
the denominator has to be a square. This reduces the time complexity
tremendously.

\subsubsection{Odd degree general}

In general, when $f$ has odd degree, the denominator has to be `almost'
a square: it must be a square times a (squarefree) divisor of the leading
coefficient, and there are further restrictions on the parity of the
valuation at~$p$, for $p$ dividing the leading coefficient, when this
valuation is sufficiently large. This gives the same type of time complexity
as in the monic case, but with a larger constant.

\subsubsection{Even degree} \label{sec:elim}

Let $l$ be the leading coefficient. If $l$ is not a square, we can exclude
denominators divisible by a prime~$p$ such that $l$ is a non-square
mod~$p$: for $p \mid b$ we have $F(a,b) \equiv l a^n \bmod p$ with $p \nmid
a$, which is a non-square as well. In addition, we use the necessary
condition that the Jacobi symbol $\left(\frac{l}{b'}\right)$ must be $+1$,
where $b'$ is the part of $b$ coprime to~$2l$. Since version~2.3 that symbol
is not computed by a binary algorithm for each denominator but as a product
of Legendre symbols: for $l = \pm 2^v \prod_i q_i^{e_i}$, with $\epsilon = 1$
if $l < 0$ and $0$ otherwise, it is
$\left(\frac{-1}{b'}\right)^{\epsilon} \left(\frac{2}{b'}\right)^{v}
\prod_i \left(\frac{q_i}{b'}\right)^{e_i}$, where by quadratic reciprocity
each factor is the Legendre symbol of $b' \bmod q_i$ modulo~$q_i$ times a
sign that depends on $b' \bmod 4$, and the first two factors depend on $b'
\bmod 8$. The program tabulates the non-squares modulo each~$q_i$ once per
curve, so that the test costs one multiplication and one table look-up per
prime factor of~$l$; when $l$ has a prime factor above $1021$, or more odd
prime factors than the tables allow for (sixteen, or 8~KB), or the
denominators exceed $2^{32}$, it falls back on the binary algorithm.

When $p$ divides~$l$, the congruence says nothing, since $F(a,b)$ is then
$0$ mod~$p$; the valuation does. Write $v_p(b) = m \ge 1$ and let $w_j$ be
the valuation at~$p$ of the coefficient of~$b^j$ in~$F$, so that the term
with~$b^j$ has valuation $w_j + jm$. If one term has strictly smaller
valuation than all the others, it determines $v_p(F(a,b))$ and the residue
of $F(a,b)/p^{v}$ mod~$p$. If that valuation is odd, no denominator with
$v_p(b) = m$ can occur. If it is even and $j$ is even, then $a^{n-j}$ and
$(b/p^m)^j$ are both squares mod~$p$, so none can occur when the unit part
of that coefficient is a non-square mod~$p$. (For odd~$j$ the factor
$(b/p^m)^j$ takes non-square values as well as square ones, and nothing
follows.) When two terms tie for the smallest valuation they may cancel, and
nothing follows either; indeed that is what happens at the points of such a
curve. For large~$m$ the leading term wins by itself, so when $v_p(l)$ is
odd, or even with a non-square unit part, every sufficiently large valuation
is excluded: the denominator cannot be divisible by a certain power of~$p$.
The commonest cases are $v_p(l) = 1$ with $p \mid l_{n-1}$, which excludes
$p$ itself, and $v_p(l) = 1$ with $p \nmid l_{n-1}$, where the two top terms
tie at $m = 1$ and $p^2$ is excluded. Below such a tail there may be isolated
excluded valuations, and there may be some when $l$ is a square: $l = 9u$
with $3 \nmid l_{n-1}$ excludes $v_3(b) = 1$ whatever $u$ is.

While computing the sieving information, we make a list of these primes,
each with the set of valuations it excludes, which we then use to eliminate
denominators: a prime that is excluded altogether is tested for with bit
arrays in the manner of the sieve, the others by computing the valuation of
the denominator. The bit arrays cover every excluded prime up to the square
root of the height bound (or up to~$1021$), building the patterns for the
primes beyond the compiled sieving primes as needed: past the square root a
denominator can have at most one excluded prime left, and the Jacobi symbol
sees that one, whereas it lets an even number of them through.

Up to version~2.2.3 the valuation test never ran. It sat behind a guard that
asked for $l$ to be a non-square modulo~$p$, and $l \equiv 0$ is a square.
That was fortunate, because it was also wrong: in working out from which
power of~$p$ on the leading term decides the valuation, it left out every
coefficient that is a unit at~$p$, which are the ones that decide it.

\subsection{Reversing the polynomial}

The exclusion of denominators described above is more or less effective,
depending on the situation; it is the better the earlier the case was
described. Since the height of the points does not change under the
transformation $(x,y) \to (1/x, y/x^{\lceil n/2 \rceil})$, we can as well search on the
reversed polynomial $F(z,x)$. This will result in a speed-up if the
reversed polynomial belongs to a `better' class than the original.

\subsection{Some general remarks}

Modern processors are very fast when doing basic things like moving
data around between registers or simple arithmetic operations (addition,
subtraction, shift, \dots, comparison). Even memory access can be fast
when there is no cache miss. On the other hand, integer division is a
slow process. It turned out that quite some improvement of the performance
was possible by removing as many instances of integer division operations
as reasonably possible.

For example, the gcd that tests a candidate for lowest terms is a binary
gcd without a data-dependent branch; the Jacobi symbol test on the
denominators is evaluated as a product of Legendre symbols read from tables
built once per curve (Section~\ref{sec:elim}); the residue classes of~$b$
modulo the various sieving primes are computed by a multiplication with a
precomputed reciprocal rather than by a division; and most of the testing of
`forbidden divisors' of~$b$ is done with bit arrays similar to those used
for sieving the numerators, a word of 64 denominators at a time.

\subsection{Change log}

\textbf{Version 2.0} was released January 9, 2008.

\textbf{Version 2.0.1} was released July 7, 2008. It fixes a bug that prevented the
'\texttt{-1}' option to work properly.

\textbf{Version 2.1} was released March 9, 2009. It makes use of the SSE instructions,
so that the sieve can work on 128~bits in parallel (instead of on 64 or~32).
On my laptop (with an Intel Core2 processor), \texttt{make test} runs about
25\% faster than before.

\textbf{Version 2.1.1} was released April 14, 2009. It fixes a bug that in some
cases prevented an early abort when \texttt{*quit} was set by the callback function.
Thanks to \textbf{Robert Miller} for the bug report and the fix. In addition, it is now
checked that \verb+__WORDSIZE == 64+ before SSE instructions are used
(in \texttt{rp-private.h}; removed again in version~2.3, which assumes a
64-bit \texttt{long} throughout), since the program then assumes that a bit
array consists of two \texttt{unsigned long}s. In this context, \verb+__SSE2__+ has been replaced
by a new macro \verb+USE_SSE+. A further fix eliminates unnecessary copying
of sieving information when SSE instructions are not used (introduced in version~2.1).
This was almost always harmless, but could have resulted in memory corruption
in the extreme case that all the information had to be computed for all the
primes.

\textbf{Version 2.1.2} was released May 27, 2009. It fixes some memory leaks.
Thanks to \textbf{Robert Miller}.

\textbf{Version 2.1.3} was released September 21, 2009. The library function
\verb+find_points+ should now work without any bound on the degree of the
polynomial. The degree bound for the \texttt{ratpoints} executable is now set
to~100 by default (specified by \verb+RATPOINTS_MAX_DEGREE+, used to be~10).

A bug in \texttt{rptest.c} (pointed out to me by \textbf{Giovanni Mascellani}
and \textbf{Randall Rathbun}) that could lead to a segmentation fault was fixed on
March 10, 2011.

\textbf{Version 2.2} was released January 9, 2022. The main change is that
\texttt{ratpoints} can now also work with 256-bit (and, in principle, 512-bit)
registers when the CPU has the relevant capabilities. Thanks to
\textbf{Bill Allombert} for doing a first implementation of this. In addition,
the code has been cleaned up to some extent (e.g., the variants according
to register sizes are now completely dealt with through macros defined
in \texttt{rp-private.h}), and some more comments have been added.
There are also some further optimizations; for example, the first sieving
phase now uses many of the SSE/AVX registers in parallel, and on some
architectures, the code uses a faster implementation for the test whether
a bit array is zero. On my current laptop, the new code (using 256-bit words)
is about twice as fast as the old one (2.1.3 using 128-bit words).

\textbf{Version 2.2.1} was released January 18, 2022. This fixes a small
bug that was introduced in the previous version, which could lead to
the program missing points at the upper end of the search interval
for the numerators. Thanks to \textbf{Bill Allombert}.

\textbf{Version 2.2.2} was released May 15, 2023. This fixes a small
bug that led to the polynomials being reversed incorrectly when the
degree is odd. Thanks to \textbf{Nicolas Mascot} and \textbf{Bill Allombert}.

\textbf{Version 2.2.3} was released September 6, 2026. It fixes the code
path using 512-bit registers, which had never been exercised. Two problems
were found. The first one would have made
the program crash: the array \texttt{sieves0} of precomputed sieving
information is declared as an array of \texttt{unsigned long}, but it is
read through pointers to \texttt{ratpoints\_bit\_array}, and \texttt{gcc}
aligns such a global array to at most 32~bytes, whereas a 512-bit
bit-array needs 64. The aligned 512-bit load that \texttt{gcc} generates
in the second sieving phase would then fault; whether it did so depended
on where the linker happened to place the array. The array is now given
the alignment of \texttt{ratpoints\_bit\_array} explicitly, in every
variant. The second problem was that the \texttt{TEST} macro, which is
used in the innermost loop of the second phase, still used the generic
fall-back version; \texttt{gcc} compiles this to a spill of the whole
64-byte vector to the stack followed by reading it back word by word,
which defeats store-to-load forwarding. It now uses \texttt{vptestmq},
in analogy with the 256-bit version.

Since I have no CPU with AVX512F available, this was tested by compiling
with \verb+-DUSE_AVX512+ but without \verb+-mavx512f+, so that
\texttt{gcc} expresses the 64-byte vectors through 256-bit operations.
This exercises the whole code path (in particular \verb+RBA_PACK == 8+
and the two mask macros); \texttt{make test1} and \texttt{make test2}
both reproduce the reference output exactly, as do 400 randomly generated
comparisons against the 256-bit version. What this cannot check is the
genuine 512-bit instructions. In the meantime, \textbf{Drew Sutherland} has
run the two variants against each other on an AMD Ryzen~9 9950X3D (Zen~5
architecture), for a curve with many rational points and a height bound
of~$10^6$; the 512-bit version was faster by 13 to 14\%.

\textbf{Version 2.3} is in preparation. Its theme is that the program should
work out for itself what it used to be told.

The number of primes used in the two sieving stages, \texttt{sp1}
and~\texttt{sp2}, is no longer a fixed pair of defaults but is chosen from the
curve, as described in Section~\ref{sec:spchoice}; that is worth about $14\%$
on the random curves of \texttt{make test1} and about $36\%$ on the point-rich
curves of the new \texttt{make test1many}. The choice rests on two constants
that depend on the machine rather than on the curve. They can be set with the
new \verb+-r+ and \verb+-R+ options, and \texttt{make tune} will measure them
for your machine and write them to \texttt{tuning.mk}
(Section~\ref{sec:tuning}).

A curve with very many rational points makes $f$ a square modulo every residue
for the smallest primes, so those primes say nothing and are discarded, and
there can be too few left for that choice to have the primes it wants. The
program now looks at further primes when, and only when, that happens; the
default prime bound has been raised from $127$ to~$251$ to give it something to
reach for. On the curves this applies to it is worth between $1.3$ and $2$
times at a height bound of~$2\cdot10^5$, and no curve that does not need it
pays anything for it.

The second sieving stage now steps over the empty bit arrays in a tight loop of
its own, and finds the surviving numerators with a count of trailing zeros
rather than by testing every bit position up to the highest one set.

Two things that had never been right in the build are fixed. Changing
\texttt{CCFLAGS1} or \texttt{PRIME\_SIZE}, or retuning, now rebuilds what
depends on it: a stamp file records the flags, so that a change of register
width no longer leaves some object files at the old one, and so that
\texttt{make tune} has an effect at all. And \texttt{CCFLAGS128} now selects
\verb+USE_AVX128+ rather than \verb+USE_SSE+, which is no less portable and
is faster, because its test for an empty register does not extract both halves
into general registers.

\verb+USE_LONG_IN_PHASE_2+ has changed meaning. It used to run the whole
second stage on words; it now keeps the scan for survivors at the register
width and sieves only the words of a surviving bit array that are themselves
non-zero. It is still not worth switching on, and it is now clear why not:
the number of \texttt{AND} steps is the same either way, because the other
words of a surviving bit array are already zero, and a narrow and a wide read
of the same table come from one cache line. While this was being done, a
latent bug came to light in the \verb+USE_SSE+ variant, where the macro
extracting word~$i$ of a bit array ignored~$i$ and always returned word~$1$;
that was harmless as long as the one caller wanted nothing else.

The memory reserved for the sieving tables now follows the number of primes
that may actually be used, rather than the whole compiled-in table. This
matters because of the raised prime bound: at the largest supported bound the
old reservation would have been 1.7~GB.

Building the sieve tables is faster by a factor of nearly four, which is worth
13\% of \texttt{make test1} and little on a run with a large height bound,
where the tables were never a large part of the work. Two things were wrong
with the loop that turns the table of squares into a bit pattern. It tested
that table with a branch, and which residues are squares has no pattern, so
about half of those branches were mispredicted; shifting the value into place
has nothing to miss. What is left is then a chain --- one row could not start
until the residue for the row before had been reduced --- and four rows are now
taken at a time, with four residues stepping by $4b^{-1} \bmod p$, so that the
four chains overlap. \verb+-DRP_INIT_BRANCH+ and \verb+-DRP_INIT_ONEWAY+ put
the old forms back, should you want to measure this on your own machine.

There are two further test suites, \texttt{make testhigh} and \texttt{make
testhighmany}, which run the two regimes at a height bound of $2\cdot10^5$
instead of $16383$. The point of them is that the old bound is short enough
for a fifth of the run to be table construction rather than sieving, so it
measures a change to the sieving loops through a filter; at the new one that
share is half a per cent. \texttt{make tunehigh} tunes the constants on
them. The test programs also no longer pass the default number of primes
explicitly, which had quietly stopped them from ever using the new rule above.

There is a third sieving stage. The two that were there work on bit arrays,
which is the right shape while a whole array's worth of numerators is still
in play; once the second stage has finished there are only a handful left per
denominator, and each of them went straight to the exact test. The new stage
tests them against further primes first, one numerator at a time, computing
where to look in the table of admissible residues rather than reading a bit
out of a sieving table. So it builds no table, which is what makes it worth
doing at that point --- a table costs work proportional to the prime for
every denominator, and by then there is about one survivor per denominator to
show for it --- and it also lets the stage use primes far up the list, which
matters most for the curves whose small primes say nothing. How many it uses
is decided per curve; on a random curve at the default height bound it uses
none, which is right. What it is worth is $1.9\%$ of \texttt{make test1many}
and $1.6\%$ of \texttt{make testhigh}, and $6.7\%$ of
\texttt{make testhighmany} --- the point-rich curves at a large height bound,
which is where the exact test had grown to $8\%$ of the run.

Reducing a word number modulo a sieving prime, which is done for every prime
at the start of every stretch of the first sieving stage and is therefore one
of the busiest operations in the program, no longer subtracts multiples of the
prime one at a time. It multiplies by a precomputed reciprocal instead, which
is what the third stage above already does. The old form was written to avoid
divisions and did avoid them, but the chain of conditional subtractions that
replaced them cost some twenty instructions on every call, against four for a
multiplication and a shift. Measured against each other inside one binary,
so that the comparison is not at the mercy of where the compiler happens to
place the surrounding loops, the new form is 2 to 3\% faster on all four test
sets; comparing whole builds instead reproduces that on the two sets at the
default height bound, and at the large one the difference between two code
layouts of the same program exceeds it. The measurement is worth stating the
other way round as well: the hardware divider was busy for at most $0.6\%$ of
the run to begin with, so the divisions were never what cost anything.

While that was being done, an undefined shift came to light in the
computation of the intervals where the polynomial is positive: an interval
endpoint was scaled by dividing by \verb+1 << del+ with \texttt{del} the
bisection depth, which is bounded by the \verb+-S+ setting and so may be as
large as $62$. It is now scaled by \texttt{ldexp}, which is exact at every
depth and needs no division.

The rule that chooses the primes has been taken further, and in three
directions at once, all of which come from noticing that the old rule valued
every prime the same and assumed every run was infinitely long. A prime's
sieving table is built at most $p$ times however many numerators there are,
and its entry in the list of reduced denominators is computed once per
denominator, so per word of numerators both fall like $1/U$ with the length
of the run. The program now works out $U$, and the number of denominators,
before it looks at a single prime. With that it (i)~scales the offset from
\texttt{sp1} to \texttt{sp2} by the length of the run, which is what lets one
tuning serve every height bound and is worth $12\%$ of \texttt{make
testhighmany}; (ii)~ranks the primes by what they cost divided by what they
say rather than by what they say alone, so that at equal quality it prefers
a smaller prime, whose table is cheaper; and (iii)~counts what the sieve is
actually doing --- the bit arrays that survive the first stage, the survivors
of the second and the ones that pass the test for common factors --- and
revises the numbers as the run goes on. The last of these is what finally
gets at the one thing the prediction could never see: the non-reduced
representations $(ka, kb)$ of a rational point pass every prime test, so a
floor of survivors outlives any amount of sieving. Sections~\ref{sec:spchoice}
and~\ref{sec:adapt} have the details, and \verb+-U 0+, \verb+-C 0+ and
\verb+-A 0+ switch the three off individually.

Together they are worth $5.6\%$ of \texttt{make test1}, $11.9\%$ of
\texttt{make testhighmany} and $13.7\%$ of the new \texttt{make testdegrees},
and cost half a per cent on the random curves at the large height bound,
which want one or two fewer second-stage primes than the pair of constants
gives them.

The loop in \texttt{sieving\_info} that evaluates the polynomial at every
residue modulo every prime now reduces by multiplying rather than by
dividing, on a fixed schedule rather than when a test says it must. As a
side effect this repairs something that raising \texttt{PRIME\_SIZE} from~7
to~8 had broken: it moved the boundary of the division-free path from
degree~8 down to degree~7, so that a genus~3 curve with an even model was
taking a conditional division inside every step of the evaluation. Nothing
had noticed, because every test curve in the package has degree~6.

The two constants that say how many primes the third stage should use are
fractions of what one exact test costs, and that is a property of the curve,
not of the machine alone: the test evaluates a binary form of the given
degree and takes an integer square root, so it costs about 15~cycles per
degree plus a step for every further limb the root needs, and over the range
from degree~3 with small coefficients to degree~14 with 400-bit ones it
varies by a factor of ten. The program now estimates it from the degree, the
largest coefficient and the height bound before it looks at a prime
(Section~\ref{sec:sp3choice}), and \verb+-W+ overrides the estimate.
Measured against the old behaviour, which was to assume that every curve
costs what a degree-6 curve with small coefficients costs, this changes no
running time by as much as half a per cent on any of the test sets --- the
rule it feeds sits at its own optimum, so a better constant moves the
boundary by about one prime, and the prime at the boundary is by definition
worth about what it costs. What it buys is that the constants mean what they
are documented to mean for curves outside the ones they were tuned on.

The first sieving phase no longer reads the bit arrays it is about to
overwrite. Every bit array starts from the same pattern of admissible
numerators modulo~16, and that pattern used to be written into all of them on
a pass of its own before the phase began; the first prime now ANDs it in as
it sieves, so the pass is gone and with it one store and one load for every
bit array --- 1.7\,$\cdot\,10^{10}$ of them in \texttt{make testhigh}. The
two boundary words are cleared after the phase instead of before it, which
comes to the same thing because \texttt{AND} is commutative. It is
worth $5.6\%$ of \texttt{make testhigh}, $5.3\%$ of \texttt{make
testdegrees} at a height bound of $2\cdot10^5$, $3.4\%$ of \texttt{make
testhighmany} and $2.6\%$ of \texttt{make test1}.

Nor does the first phase sieve bit arrays that are then thrown away. It works
on \texttt{RATPOINTS\_CHUNK} $= 16$ bit arrays at a time, one per vector
register, and the range of every denominator and interval used to be padded
up to a multiple of that: the padding was sieved with every prime of the
phase, walked by the scan of the second phase and then zeroed, and it came to
a seventh of all the bit arrays swept in \texttt{make test1}, a quarter of
them at a height bound of 4000 and four fifths at 1000, where the whole
interval of a denominator is three bit arrays long. Now the chunk loop stops
at the last whole chunk, and the bit arrays left over --- fewer than
sixteen --- are sieved in legs of 8, 4, 2 and 1 registers, one leg for each
set bit of their number. A leg is cheaper than a chunk per prime, and not
only for being narrower: it reads its rows from the pointers the chunk loop
left, up to fifteen rows past them, which is what the wrap-around copies at
the end of every table are there for, so it neither wraps the pointers
around nor stores them back --- nothing reads them after the phase, and the
next call recomputes them. Worth $10\%$ of \texttt{make test1} ($9.5$ to
$11.4\%$ at three placements of the code), $6.5\%$ of \texttt{make
test1many}, $2\%$ of \texttt{make testhigh}, between nothing and $2.5\%$
of \texttt{make testhighmany}, and a fifth of the run at the height bounds
1000 and 4000.

The test on the denominators at a prime~$p$ dividing the leading coefficient
now runs, and is right. It had been unreachable since it was written, behind
a guard that asked for the leading coefficient to be a non-square modulo~$p$
--- which zero never is --- and it was unsound as well: in working out from
which power of~$p$ on the leading term of $F(a,b)$ decides the valuation, it
skipped every coefficient that is a unit at~$p$, as though their valuation
were infinite rather than zero, and it tested the unit part of the leading
coefficient for being a square with its sign dropped. Enabled as it stood, it
lost 23 points on 12 of the thousand test curves. The test is now derived
from the Newton polygon of~$F$ at~$p$ (Section~\ref{sec:elim}): for each
valuation a denominator can have at~$p$ it asks whether a single term of
$F(a,b)$ sets the valuation, and excludes the ones where it does with an
odd valuation or a non-square unit. It applies to a square leading
coefficient too, which used to switch every test on the denominators off.
It excludes something on 499 of the thousand random curves of \texttt{make
test1}, whose coefficients are small enough for 3, 5 and 7 to divide the
leading one often, and is worth $8.3\%$ there, $5.3\%$ of \texttt{make testdegrees} and
$2.4\%$ of \texttt{make test1many}, $8.8\%$ of \texttt{make testhigh}
and $2.1\%$ of \texttt{make testhighmany}. That is more than the share
of denominators it excludes, because a prime that divides the leading
coefficient passes every numerator of a denominator it divides, so those
denominators reach the later stages with several times the usual survivors.

A review of the code in September 2026 turned up five bugs, which are fixed
here and, where they are older than this version, in 2.2.4 as well. A
polynomial of degree~1 --- given as such, or reached from
$y^2 = c_2x^2 + c_1x$ by the reversal, or from a zero leading coefficient
--- made the Sturm sequence code read past the end of its arrays and crash.
When the positivity region of~$f$ did not meet the search domain, the
program returned before it had looked at the points at infinity, so that
\texttt{ratpoints '-1000000000 0 0 1' 45} printed nothing where
\verb+-s+ printed $(1:0:0)$, and after a reversal the lost point is an
affine point of the given curve. When no residue class of the denominator
modulo~16 admitted any numerator, the array of numerator patterns was read
before it had been filled; by default the answer was still right, but the
choice of parameters was not reproducible, and with \verb+-j -F 0+ the
program printed nondeterministic points outside the height bound after a
run several thousand times longer than necessary. When no {\em odd}
denominator admitted a numerator but even ones did, the estimate of the
run length fell to its floor, which switched the second and third sieving
stages and the adaptive correction off for those curves (41 of the 900
random test curves that sieve; the points were right, the run slow). And
several input fields of \texttt{ratpoints\_args} were used as scratch
space --- the number of forbidden divisors found was written into
\texttt{max\_forbidden}, the intersected search intervals into
\texttt{num\_inter}, the chosen numbers of primes into \texttt{sp1} and
\texttt{sp2}, normalised values into several others --- so that a program
which fills the structure once and then calls \texttt{find\_points\_work}
for curve after curve, as Section~\ref{sec:library} suggests, had the
forbidden-divisor test switched off after the first curve without forbidden
divisors (a factor of two in running time), the first curve's choice of
primes imposed on all the others, and its positivity region intersected
into their search domains, which loses points. The input fields now come
back as they went in, except \texttt{num\_inter} and \texttt{domain}, which
describe the region searched and are documented as such, and
\texttt{sp1\_used}, \texttt{sp2\_used} and \texttt{sp3\_used} report what
was used; \texttt{make test1once} runs the test curves that way.
\texttt{make test3} runs the invocations of \texttt{test3.sh}, one or more
for each of these bugs, against \texttt{testbase3}.
Also fixed: four left shifts of negative values (undefined in C, harmless
with gcc so far), two arrays that could be declared with length zero under
\verb+-n 0 -N 0+, an overflow in the enumeration of the divisors of the
leading coefficient for denominator bounds above about $4\cdot10^{15}$, and
two constants in \texttt{ratpoints.h} that could not be overridden from the
compiler command line.

The code now assumes that \texttt{long} has 64~bits and refuses to compile
otherwise, so the 32-bit architectures are no longer supported; what used
to allow for a 32-bit word has been removed. Version~2.2.4 is the last one
that accommodates a 32-bit \texttt{long}.

Two loops of the second sieving phase have been rewritten. The scan that
steps over the empty bit arrays after the first phase used to carry a bound
test and two counters, nine instructions per bit array; it now runs into a
non-zero sentinel written just past the range, and is a load, a test, a
branch and the pointer step. And the phase used to reach the table row of a surviving bit array
from a pointer set up at the start of every call, subtracting the prime until
the pointer was back inside the table --- one to three data-dependent, mostly
mispredicted branches for every \texttt{AND}, and one such set-up per prime
and call whether or not anything had survived. It now computes the row from
the word number, by the reduction with a precomputed reciprocal that the
first phase and the third stage already use, with a multiple of the prime
folded into the offset so that the value reduced is never negative. Together
the two are worth $3\%$ of \texttt{make test1} and of \texttt{make
test1many}, and $7.5\%$ of \texttt{make testhigh} and of \texttt{make
testhighmany} at the default code placement, between $4.5$ and $7.5\%$ at
others: the first by taking instructions out of the busiest short loop of
the program, the second by taking out branch mispredictions, which fall by a
fifth to nearly a third at the large height bound.

\texttt{make tune} also measures a fourth constant now, what building one
row of a sieving table costs, which is what decides how strongly the program
prefers a smaller prime to a larger one at a small height bound. It could be
set with \verb+-C+ since the rule was introduced, but nothing had measured
it, and of the four constants it is the one most likely to differ between
machines, being the ratio of two different kinds of work.

The loop over the denominators, and what is done once per denominator, have
been reworked, six changes together. The Jacobi symbol test on the
denominators no longer computes a symbol for each denominator by a binary
algorithm but evaluates it as a product of Legendre symbols read from tables
built once per curve (Section~\ref{sec:elim}); that computation was $12\%$ of
\texttt{make test1} and a third of its branch mispredictions. The tests that
depend on $b \bmod 64$ alone are applied to a word of 64 denominators at
once, and the loop visits only the denominators that pass them. The residues
of~$b$ modulo the sieving primes are computed by a multiplication with a
precomputed reciprocal instead of being stepped from the previous denominator
by conditional subtractions, which were mispredicted a quarter of the time.
The gcd that tests a candidate for lowest terms has no data-dependent branch
left. The bit arrays that test for forbidden divisors cover every excluded
prime up to the square root of the height bound rather than only the compiled
sieving primes, and the default of \verb+-F+ is $64$ instead of $30$. And what
the third stage needs for a denominator is set up when the first candidate
reaches it rather than for every denominator. Measured step by step against
the version before, the six are worth $14$, $2$, $7$, $2$, $0$ and $2\%$ of
\texttt{make test1}; together they are $24\%$ of \texttt{make test1}, $4\%$
of \texttt{make test1many}, $6\%$ of \texttt{make testhigh} and $2\%$ of
\texttt{make testhighmany} at the default code placement, and within a
point of that at two other placements, and two thirds of the branch
mispredictions of \texttt{make test1} are gone.

Finally, the \texttt{ratpoints\_args} structure has gained fields, so
programs using the library must be recompiled against the new header.

%=========================================================================

\begin{thebibliography}{gmp}
\frenchspacing

\bibitem[Coh]{Cohen}
  H. Cohen, \textit{A course in computational algebraic number theory,}
  Springer GTM~\textbf{138}, second corr. printing, 1995.

\bibitem[gmp]{gmp}
  The GNU Multiple Precision Arithmetic Library, \url{http://gmplib.org/}

\bibitem[rat]{ratpoints}
  The \texttt{ratpoints-\rpversion} package. Available at \\
  \url{http://www.mathe2.uni-bayreuth.de/stoll/programs/ratpoints-\rpversion.tar.gz}

\end{thebibliography}

\end{document}
